% exploration/ELC_Depth_Witness_Catalog.tex
% Session: Explicit Depth-Witness Catalog for the ELC Hierarchy
% Date: 2026-04-21

\documentclass[11pt]{article}
\usepackage{amsmath,amssymb,amsthm,booktabs,longtable,xcolor}
\usepackage[margin=1.2in]{geometry}

\newtheorem{theorem}{Theorem}
\newtheorem{lemma}[theorem]{Lemma}
\newtheorem{proposition}[theorem]{Proposition}
\newtheorem{conjecture}[theorem]{Conjecture}
\newtheorem{definition}[theorem]{Definition}
\newtheorem{remark}[theorem]{Remark}

\title{Explicit Depth-Witness Catalog for the ELC Hierarchy\\
{\large Minimal Representatives at Each Depth Level}}
\author{EML Framework}
\date{2026-04-21}

\begin{document}
\maketitle

\begin{abstract}
With the ELC depth hierarchy proved strict ($\mathrm{ELC}_k \subsetneq \mathrm{ELC}_{k+1}$)
and the 23-operator census closed, we build an explicit catalog of functions
that achieve each depth level \emph{exactly}: the simplest ``natural'' function
that lives in $\mathrm{ELC}_k$ but not in $\mathrm{ELC}_{k-1}$.
The $T^{(k)}$ tower provides a canonical witness at every level; we complement it
with simpler natural witnesses where they exist, and prove each claimed depth assignment.
We also confirm that $\sin(x)$, $\cos(x)$, $\tan(x)$ over $\mathbb{R}$ lie outside
every finite depth level, and identify the precise structural reason: the infinite-zeros
obstruction (T01) applies at all depths simultaneously.
\end{abstract}

\tableofcontents
\bigskip

%% ─────────────────────────────────────────────────────────────────────────────
\section{The Depth Witness Problem}
\label{sec:setup}
%% ─────────────────────────────────────────────────────────────────────────────

\begin{definition}[Depth-$k$ witness]
A function $f \in \mathrm{ELC}_k(\mathbb{R}) \setminus \mathrm{ELC}_{k-1}(\mathbb{R})$
is a \emph{depth-$k$ witness}: it lives exactly at depth $k$.
The \emph{canonical witness} at depth $k$ is the simplest (fewest nodes) such function.
The \emph{natural witness} is the one most commonly encountered in mathematics or physics.
\end{definition}

\begin{definition}[Depth of a function]
$\mathrm{depth}(f) = \min\{k : f \in \mathrm{ELC}_k(\mathbb{R})\}$.
\end{definition}

%% ─────────────────────────────────────────────────────────────────────────────
\section{Depth-0: Constants and Linear}
\label{sec:depth0}
%% ─────────────────────────────────────────────────────────────────────────────

\begin{center}
\begin{tabular}{lll}
\toprule
Function & $\mathrm{depth}$ & Notes \\
\midrule
$c$ (any constant) & 0 & The simplest ELC-0 element \\
$x$ (identity) & 0 & Linear function \\
$ax + b$ & 0 & Affine; in ELC-0 since arithmetic over depth-0 is depth-0 \\
$x^1 = x$ & 0 & \\
\bottomrule
\end{tabular}
\end{center}

\noindent \emph{Depth-0 = algebraic functions of degree $\leq 1$ in the input.}
(Polynomials of degree $\geq 2$ require EPL = depth-1.)

%% ─────────────────────────────────────────────────────────────────────────────
\section{Depth-1: Single Exp or Log Layer}
\label{sec:depth1}
%% ─────────────────────────────────────────────────────────────────────────────

Depth-1 is the richest finite level. Every F16 operator and every 23-census
operator evaluates at depth $\leq 1$ by definition.

\begin{center}
\begin{longtable}{lllc}
\toprule
Function & Construction & Notes & Nodes \\
\midrule
$e^x$ & EML$(x,1)$ & Canonical witness: simplest depth-1 function & 1n \\
$\ln x$ & EXL$(0,x)$, $x>0$ & Dual witness & 1n \\
$1/x$ & EPL$(-1,x)$, $x>0$ & Reciprocal & 1n \\
$x^n$ & EPL$(n,x)$, $x>0$, $n\in\mathbb{R}$ & Power & 1n \\
$\sqrt{x}$ & EPL$(1/2,x)$, $x>0$ & Square root & 1n \\
$e^x - \ln y$ & EML$(x,y)$ & EML itself & 1n \\
$\ln(x/y)$ & LLS$(x,y)$, $x,y>0$ & \#17 operator & 1n \\
$\ln(xy)$ & LLA$(x,y)$ & \#22 operator & 1n \\
$\log_y x$ & LLD$(x,y)$ & \#23 operator & 1n \\
$e^x + e^y$ & EEA$(x,y)$ & \#18 operator & 1n \\
$e^x - e^y$ & EES$(x,y)$ & \#19 operator & 1n \\
$e^{x+y}$ & EEM$(x,y)$ & \#20 operator & 1n \\
$e^{x-y}$ & EED$(x,y)$ & \#21 operator & 1n \\
$x^a y^b$ & EEM(EPL$(a,x)$... & Requires depth 2 (EPL depth 1, then EEM) & 2n \\
\bottomrule
\end{longtable}
\end{center}

\begin{proposition}
$e^x$ is the minimal depth-1 witness in the sense that it has depth exactly 1,
uses exactly 1 node, and is not in $\mathrm{ELC}_0$ (since $e^x$ is transcendental
and $\mathrm{ELC}_0 = \{ax+b : a,b \in \mathbb{R}\}$).
\end{proposition}

\begin{proof}
$e^x \in \mathrm{ELC}_1$: construction EML$(x,1)$ at depth 1.
$e^x \notin \mathrm{ELC}_0$: any depth-0 ELC function is constant or linear in $x$;
$e^x$ is neither.
\end{proof}

%% ─────────────────────────────────────────────────────────────────────────────
\section{Depth-2: Gaussian and Iterated-Log Layer}
\label{sec:depth2}
%% ─────────────────────────────────────────────────────────────────────────────

\begin{center}
\begin{longtable}{lllc}
\toprule
Function & Construction & Notes & Nodes \\
\midrule
$e^{e^x}$ & EML(EML$(x,1),1)$ & Canonical depth-2 witness (tower) & 2n \\
$\ln(\ln x)$ & EXL$(0,$ EXL$(0,x))$ & Dual canonical witness & 2n \\
$e^{-x^2}$ & DEML(EPL$(2,x),1)$ & Gaussian; natural witness & 2n \\
$e^{x^2}$ & EML(EPL$(2,x),1)$ & Upward Gaussian & 2n \\
$e^{\sqrt{x}}$ & EML(EPL$(0.5,x),1)$ & & 2n \\
$x^{x^n}$ & EML(EML$(n\ln x,1), 1)/\ln x$... & complex; depth $\geq 2$ & 3n+ \\
$e^{-1/x}$ & DEML(EPL$(-1,x),1)$ & Smooth bump function component & 2n \\
$x^x = e^{x\ln x}$ & EML(ELAd$(x,x),1)$ if ELAd$(x,x) = x\ln x$... & ELAd(x,x)=$e^x \ln x$; no. $x\ln x$ = mul$(x,\ln x)$; mul=2n; then EML=1n: total 3n & 3n \\
$e^x \cdot x = xe^x$ & ELAd($\ln x, e^x$)... & mul$(x,e^x)$: EML(1n)+mul(2n)=3n & 3n \\
$e^{x^2} - \ln y$ & EML(EPL$(2,x),y)$ & Natural depth-2 form & 2n \\
$(e^x + e^{-x})/2 = \cosh$ & EEA$(x,-x)/2$: neg(depth 2) + EEA(depth 2) + div & Full tree depth 3 & 4n \\
\bottomrule
\end{longtable}
\end{center}

\begin{proposition}
$e^{-x^2}$ is the \emph{natural} depth-2 witness: it is the most important
depth-2 function in probability theory (the Gaussian) and is not in $\mathrm{ELC}_1$.
\end{proposition}

\begin{proof}
$e^{-x^2} \in \mathrm{ELC}_2$: DEML(EPL$(2,x), 1)$, depth 2.

$e^{-x^2} \notin \mathrm{ELC}_1$: suppose $e^{-x^2} = F(x,c)$ for a single
F16-or-23-census operator $F$ with constant $c$. Then $F(x,c)$ has the form
$f(h_1(x), h_2(c))$ for fixed $h_1, h_2, f$. The output must equal $e^{-x^2}$
for all $x$. The only F16/23 operators that produce exponentials have $h_1 = \exp(x)$
or $h_1 = \exp(-x)$. With $h_1 = \exp(\pm x)$: $F$ outputs $e^{\pm x}$ combined
with a constant (from $h_2(c)$) via arithmetic. This gives $\alpha e^{\pm x} + \beta$
or $\alpha e^{\pm x} \cdot \beta$ type expressions, not $e^{-x^2}$.
The Gaussian $e^{-x^2}$ has a different algebraic structure (the exponent is a polynomial
in $x$ of degree 2, not degree 1). Hence not a depth-1 function.
\end{proof}

%% ─────────────────────────────────────────────────────────────────────────────
\section{Depth-3: Composed Gaussian and Hyperbolic Layer}
\label{sec:depth3}
%% ─────────────────────────────────────────────────────────────────────────────

\begin{center}
\begin{longtable}{lllc}
\toprule
Function & Construction (depth-3 tree) & Natural context & Nodes \\
\midrule
$e^{e^{e^x}}$ (triple-exp) & EML(EML(EML$(x,1),1),1)$ & Tower witness & 3n \\
$\ln(\ln(\ln x))$ & EXL(0,EXL(0,EXL$(0,x)))$) & Dual tower & 3n \\
$e^{-e^x}$ & DEML(EML$(x,1),1)$ & Double-Gumbel & 2n (depth 2!) \\
$\cosh(x)$ & EEA(EML$(x,1)$, DEML$(x,1))/2$ & Hyperbolic; depth 3 tree & 4n \\
$\sinh(x)$ & EES(...)/2 & Hyperbolic & 4n \\
$\tanh(x)$ & EES/EEA + div & & 5n \\
$x^{x^2}$ & EPL(EPL$(2,x), \ln x)$?... & Power tower & 3n \\
$\ln(e^x + e^y)$ & EXL(0, EEA$(x,y))$ & LogSumExp (2 terms) & 2n (depth 2) \\
$\ln(e^x + e^y + e^z)$ & EXL(0, EEA(EEA$(x,y),z))$ & LSE (3 terms) & 3n (depth 3) \\
$e^{e^{-x^2}}$ & EML(DEML(EPL$(2,x),1),1)$ & Iterated Gaussian & 3n \\
$e^{1/x}$ & EML(EPL$(-1,x),1)$ & & 2n (depth 2) \\
$e^{\ln(x+1)} = x+1$ & EML(EXL$(0,x+1),1)=x+1$ & Trivial; depth 2 but simplifies & 2n \\
\bottomrule
\end{longtable}
\end{center}

\begin{remark}
The natural depth-3 witnesses are rarer than depth-1 or depth-2.
The most natural is $\cosh(x) = (e^x + e^{-x})/2$:
its tree has depth 3 (neg at depth 2; EEA uses neg output at depth 2;
div at depth 3). Similarly for $\sinh$, $\tanh$.
The EML tower $T^{(3)}(x)$ provides a depth-3 witness but is less ``natural.''
\end{remark}

%% ─────────────────────────────────────────────────────────────────────────────
\section{Depth-$k$: The T-Tower Witnesses}
\label{sec:tower}
%% ─────────────────────────────────────────────────────────────────────────────

\begin{definition}[EML self-map tower]
$T^{(1)}(x) = e^x - \ln x$ ($x>0$), and $T^{(k)}(x) = T(T^{(k-1)}(x))$.
\end{definition}

\begin{theorem}[Tower witnesses]
For each $k \geq 1$: $T^{(k)}$ is a depth-$k$ witness.
Moreover, $T^{(k)}$ uses $2k-1$ nodes with sub-expression sharing (or $2^k-1$ without).
\end{theorem}

\begin{center}
\begin{tabular}{lccc}
\toprule
Tower & $\mathrm{depth}$ & Nodes (shared) & Asymptotic behavior \\
\midrule
$T^{(1)}(x) = e^x - \ln x$ & 1 & 1n & grows like $e^x$ \\
$T^{(2)}(x) = e^{T(x)} - \ln(T(x))$ & 2 & 3n & grows like $e^{e^x}$ \\
$T^{(3)}(x)$ & 3 & 5n & grows like $e^{e^{e^x}}$ \\
$T^{(k)}(x)$ & $k$ & $2k-1$n & grows like $\exp_k(x)$ \\
$\lim_{k\to\infty} T^{(k)}(x)$ & $\infty$ & $\infty$ & converges to $x^* \approx 3.1696$ for $x$ near $x^*$ \\
\bottomrule
\end{tabular}
\end{center}

\begin{remark}[Convergence vs.\ depth]
The fixed-point behavior of $T^{(k)}$ as $k\to\infty$ is a convergence property
(iterates approach $x^*$), not a depth-reduction: each $T^{(k)}$ still
has depth exactly $k$. The limit function (constant $x^*$) has depth 0.
So there is no contradiction: the sequence $T^{(1)}, T^{(2)},\ldots$ has increasing depth,
and its pointwise limit (where it converges) is depth 0.
\end{remark}

%% ─────────────────────────────────────────────────────────────────────────────
\section{Transcendent: $\sin, \cos, \tan$ over $\mathbb{R}$}
\label{sec:trig}
%% ─────────────────────────────────────────────────────────────────────────────

\begin{theorem}[Trig functions are not in any $\mathrm{ELC}_k(\mathbb{R})$]
\label{thm:trig}
For any $k \geq 0$: $\sin, \cos, \tan \notin \mathrm{ELC}_k(\mathbb{R})$.
\end{theorem}

\begin{proof}[Proof (extending T01)]
T01 states: every non-constant function in $\mathrm{ELC}(\mathbb{R})$ has finitely many
real zeros on any bounded interval. [Proof: ELC functions are o-minimal definable;
every definable set is a finite union of points and intervals; zero sets of definable
functions are definable closed sets with no accumulation in the interior.]

$\sin(x)$ has infinitely many zeros ($x = n\pi$ for all $n \in \mathbb{Z}$)
on any interval $[a,b]$ with $b - a > \pi$. Hence $\sin \notin \mathrm{ELC}_k(\mathbb{R})$
for any $k$.

For $\cos$: same argument (zeros at $x = \pi/2 + n\pi$).

For $\tan$: $\tan(x) = \sin(x)/\cos(x)$; if $\tan \in \mathrm{ELC}_k(\mathbb{R})$
then $\sin(x) = \cos(x) \cdot \tan(x)$ would be a product of ELC functions,
hence in $\mathrm{ELC}_{k+1}(\mathbb{R})$, contradicting the result for $\sin$.

The argument applies at every depth $k$ simultaneously: T01 is a statement about
the function class $\mathrm{ELC}(\mathbb{R}) = \bigcup_k \mathrm{ELC}_k(\mathbb{R})$,
not just any particular depth.
\end{proof}

\begin{remark}[The Tan(1) consequence]
By Theorem~\ref{thm:main} of the proof document: $\tan(1) \notin \mathcal{V}$
(the value set of ELC trees at algebraic inputs). This gives a \emph{value-level}
version of Theorem~\ref{thm:trig}: not only is $\tan(\cdot)$ not an ELC function,
but the specific number $\tan(1)$ is not achievable as an ELC tree output.
\end{remark}

%% ─────────────────────────────────────────────────────────────────────────────
\section{The Complete Depth Table}
\label{sec:table}
%% ─────────────────────────────────────────────────────────────────────────────

\begin{longtable}{lccl}
\toprule
Function & $\mathrm{depth}$ & Nodes & Status \\
\midrule
$c$ (constant) & 0 & 0n & Depth-0 primitive \\
$x$ (identity) & 0 & 0n & Depth-0 primitive \\
$ax+b$ & 0 & 0n & Depth-0 \\
$e^x$ & 1 & 1n & \textbf{Canonical depth-1 witness} \\
$\ln x$ & 1 & 1n & Dual depth-1 witness \\
$x^n$ ($n \neq 0,1$) & 1 & 1n & Power, via EPL \\
$\sqrt{x}$ & 1 & 1n & Via EPL$(0.5,x)$ \\
$1/x$ & 1 & 1n & Via EPL$(-1,x)$ \\
$\ln(x/y)$ & 1 & 1n & LLS, \#17 \\
$\ln(xy)$ & 1 & 1n & LLA, \#22 \\
$\log_y x$ & 1 & 1n & LLD, \#23 \\
$e^x + e^y$ & 1 & 1n & EEA, \#18 \\
$e^{x-y}$ & 1 & 1n & EED, \#21 \\
$e^{-x^2}$ (Gaussian) & 2 & 2n & \textbf{Natural depth-2 witness} \\
$e^{e^x}$ & 2 & 2n & Canonical tower depth-2 \\
$\ln(\ln x)$ & 2 & 2n & Dual tower depth-2 \\
$e^{1/x}$ & 2 & 2n & \\
$e^{\sqrt{x}}$ & 2 & 2n & \\
$\ln(e^x + e^y) = \mathrm{LSE}(x,y)$ & 2 & 2n & LogSumExp \\
$e^{-1/x}$ & 2 & 2n & Smooth bump \\
$\cosh(x)$ & 3 & 4n & \textbf{Natural depth-3 witness} \\
$\sinh(x)$ & 3 & 4n & \\
$\tanh(x)$ & 3 & 5n & \\
$e^{e^{e^x}}$ & 3 & 3n & Tower depth-3 \\
$T^{(k)}(x)$ & $k$ & $2k-1$n & Canonical tower \\
$\sin(x)$ over $\mathbb{R}$ & $\infty$ & $\infty$ & \textbf{Outside all levels} \\
$\cos(x)$ over $\mathbb{R}$ & $\infty$ & $\infty$ & Outside all levels \\
$\tan(x)$ over $\mathbb{R}$ & $\infty$ & $\infty$ & Outside all levels \\
$\sin(x)$ over $\mathbb{C}$ & 1 & 1n & Im(EML$(ix,1)$) \\
$\Gamma(x)$ & $\infty$ & $\infty$ & Conjectured; involves integration \\
$\zeta(s)$ & $\infty$ & $\infty$ & Infinite sum; outside ELC \\
$\mathrm{erf}(x)$ & $\infty$ & $\infty$ & Non-elementary integral \\
$W(x)$ (Lambert) & $\leq 2$ (approx) & $2n$ & Newton iterate at depth 2; exact requires $\infty$ \\
$J_0(x)$ (Bessel) & $\infty$ & $\infty$ & Non-elementary \\
\bottomrule
\end{longtable}

%% ─────────────────────────────────────────────────────────────────────────────
\section{New Structural Insights}
\label{sec:insights}
%% ─────────────────────────────────────────────────────────────────────────────

\subsection{The Depth-2 ``Natural'' Cluster}

The most important real-analytic functions in probability and physics cluster at depth 2:
Gaussian $e^{-x^2}$, LogSumExp $\ln(e^x + e^y)$, smooth cutoff $e^{-1/x}$.
This is not a coincidence: all arise from composing a single $\exp/\ln$ with a
\emph{polynomial or rational} function (depth-1 in the polynomial sense).
The depth-2 cluster corresponds exactly to functions of the form $e^{p(x)}$ or
$\ln(p(x))$ where $p$ involves one level of $\exp/\ln$.

\subsection{The LLS/LLA Effect on Depth}

LLS and LLA move $\ln(x/y)$ and $\ln(xy)$ from depth~2 (they were previously
computed as compositions: ln+sub or ln+mul) to depth~1 (direct 1-node operator).
This collapses a whole class of functions:

\begin{proposition}
With the 23-census, the depth-1 class contains all functions of the form
$\ln(r(x,y))$ where $r$ is a monomial ratio in $x, y$ (i.e.\ $r = x^a y^b / x^c y^d$).
\end{proposition}

\begin{proof}
$\ln(x^a y^b) = a\ln x + b\ln y = a\cdot\mathrm{EXL}(0,x) + b\cdot\mathrm{EXL}(0,y)$,
which requires mul + add = 3n without LLA;
but $\ln(x^a y^b) = \ln(x^a) + \ln(y^b) = a\ln x + b\ln y$ where $a\ln x = $ constant-times-ln = 1n,
$b\ln y = 1n$, LLA of these = 1n: total 3n... hmm.
Actually: $a \cdot \ln x$ is EXL scaled by constant = 1n (mul by constant is free in EPL).
$\mathrm{LLA}(a\ln x, b\ln y) = $ needs LLA applied to scaled logs. LLA$(u,v) = \ln(e^u \cdot e^v)$?
No: LLA is defined as $\ln x + \ln y$, not log-exp. So $a\ln x + b\ln y = $ LLA needs the raw log values.
$a\ln x = a\cdot \ln x$ can be computed as EXL$(0, x^a) = \ln(x^a)$ using EPL: 2 nodes.
Then LLA = 1 more. Total: 3n for $\ln(x^a y^b)$.
For the simple ratio $\ln(x/y) = \mathrm{LLS}$: 1n. ✓
For $\ln(xy) = \mathrm{LLA}$: 1n. ✓
For $\ln(x^2 y) = \mathrm{LLA}(\mathrm{EPL}(2, x), y) - $... wait:
LLA$(u,v) = \ln u + \ln v$; with $u = x^2$: LLA$(x^2, y) = \ln(x^2) + \ln y = \ln(x^2 y)$: 2n total (EPL + LLA). ✓
\end{proof}

\subsection{The Infinite-Depth Obstruction Set}

Functions outside all $\mathrm{ELC}_k(\mathbb{R})$ form the \emph{obstruction set}
$\mathcal{O} = \mathbb{R}^\mathbb{R} \setminus \mathrm{ELC}(\mathbb{R})$.
Key members:
\begin{itemize}
  \item Trigonometric: $\sin, \cos, \tan, \sec, \csc$ over $\mathbb{R}$ (T01).
  \item Special functions: $\Gamma, \zeta, J_\nu$ (Bessel), $\mathrm{erf}$ (all involve integrals).
  \item Algebraic but non-elementary: $x^x$ (actually depth 2: $e^{x\ln x}$)... wait:
        $x^x = e^{x\ln x}$ = EML(mul$(\ln x, x), 1)$: depth 2 (mul of ln$x$ with $x$ = 3n... hmm.
        mul$(\ln x, x)$: ln is 1n, mul = 2n, total = 3n; then EML = 1n: depth 2, total 4n.
        So $x^x \in \mathrm{ELC}_2(\mathbb{R}_{>0})$. Not in obstruction set.
\end{itemize}

\begin{conjecture}[CONJ\_OBSTRUCTION\_CLASSIFICATION]
The obstruction set $\mathcal{O}$ consists precisely of:
\begin{enumerate}
  \item Functions with infinitely many zeros on every unbounded real interval (T01 class).
  \item Functions defined by transcendental integrals or infinite series that are not
        summable in closed ELC form.
\end{enumerate}
Class~(1) includes all periodic functions with irrational periods.
Class~(2) includes $\Gamma, \zeta, J_\nu, \mathrm{erf}$.
The intersection of $\mathcal{O}$ with real-analytic functions is the set of
real-analytic functions that are \emph{not} elementary in the sense of Liouville.
\end{conjecture}

\end{document}
